\chapter{Euler-Bernoulli Beam Theory}
\label{chap:euler-bernoulli}

\section{Introduction}

Euler--Bernoulli beam theory originates in the classical theory of the
\emph{elastica}: the study of the equilibrium shapes of thin elastic rods and
strips. % TODO: Cite work if possible
Historically, the model did not arise as a formal reduction of
three-dimensional elasticity. Rather, it emerged from the observation that a
slender body resists bending through the curvature of its centreline. In its
simplest form, this idea is expressed by a local moment--curvature law,
\begin{equation*}
	M = EI \kappa,
\end{equation*}
where \(M\) is the bending moment, \(\kappa\) is the centreline curvature,
and \(EI\) is the flexural rigidity. Jacob Bernoulli was among the first to
formulate the problem of the elastic curve in this way, relating the bending
of a slender body to its curvature. Euler later developed this theory using
the calculus of variations, giving a systematic mathematical treatment of
elastic curves and their equilibrium configurations.

In modern language, Euler--Bernoulli theory may be interpreted as a
one-dimensional model for slender elastic bodies. The beam is represented by
its centreline, and deformation is described primarily by bending of this
curve. The bending energy is taken to be quadratic in curvature,
\begin{equation*}
	\mathcal{E}_b = \frac{1}{2}\int EI \kappa^2 \, ds,
\end{equation*}
so that equilibrium configurations follow from minimizing this energy subject
to applied loads and boundary conditions. For small transverse deflections
\(w(x)\), the curvature is approximated by \(\kappa \simeq w''(x)\), and the
bending moment becomes
\begin{equation*}
	M = -EI w''.
\end{equation*}
Combining this constitutive relation with force and moment balance gives the
classical Euler--Bernoulli beam equation,
\begin{equation*}
	\frac{d^2}{dx^2}\left(EI \frac{d^2 w}{dx^2}\right) = q(x),
\end{equation*}
or, for constant \(EI\),
\begin{equation*}
	EI w'''' = q(x),
\end{equation*}
where \(q(x)\) is the transverse load per unit length.

The connection with three-dimensional elasticity is therefore best understood
as a later continuum-mechanical interpretation of the same model. Under an
appropriate set of assumptions, the Euler--Bernoulli equations can be recovered
from three-dimensional elastic theory, as shown in the following section. More
advanced beam and plate theories, including modern geometrically exact beam
theories, can likewise be obtained through related reductions under different
kinematic and constitutive assumptions. However, this agreement should not
obscure the historical development of the theory: interpreting Euler--Bernoulli
beam theory primarily as a reduction of three-dimensional elasticity is, from a
historical standpoint, anachronistic.

\section{Derivation from Contiuum Elastic Theory}
\label{sec:derivation-from-continumm-elastic-theory}

\begin{figure}[t]
	\centering
	\def\svgwidth{\textwidth}
	\import{figure/vector}{beam-diagram-with-R.pdf_tex}
	\caption{Diagram of the coordinate system and orthonormal frame of a three-dimensional beam: Material centerline \(\vec r(s)\) maps \(s \in [0,L]\) to Cartesian \(\mathbb{R}^3\), and the orthonormal directors \(\{\vec t, \vec n, \vec b\}\).}
	\label{fig:coordinate-system}
\end{figure}

\begin{table}[b]
	\centering
	\begin{tabularx}{\textwidth}{>{\raggedright\arraybackslash}p{0.23\textwidth}>{\raggedright\arraybackslash}X}
		Unshearable         & Sections remain plane and orthogonal to \(\vec t\). This gives us a mapping \(\vec R = \vec r + \eta \vec n + \beta \vec b\) with shear strains \(\varepsilon_{tn}= \varepsilon_{tb}=0\) (and thus \(\tau_{t n}=\tau_{t b} = 0\)). \\
		Inextensible        & \(\varepsilon_m(s)=0\) (midline/membrane strain), equivalently \(\abs{\vec r'(s)}=1\) and axial force \(N=0\).                                                                                                                     \\
		No torsion          & Bishop frame: twist rate \(\kappa_t = \vec \kappa \cdot \vec t=0\) and torsional moment \(M_t=0\).                                                                                                                                 \\
		Isotropic bending   & \(I_{\beta\beta}=I_{\eta\eta}\equiv I\) with homogeneous \(E\).                                                                                                                                                                    \\
		No off-axis bending & Choose principal bending axes s.t. \(I_{\eta\beta}=0\)                                                                                                                                                                             \\
		Small strain        & Use \(\sigma_{tt} = E \varepsilon_{tt}\) and \(\psi = \frac12 E \varepsilon_{tt}^2\).                                                                                                                                              \\
	\end{tabularx}
	\caption{Table of assumptions used to derive the classical Euler--Bernoulli beam theory from elastic theory, used in the derivation of \ref{sec:derivation-from-continumm-elastic-theory}.}
	\label{tab:assumptions-euler-bernoulli-beam}
\end{table}

Consider a beam with a centerline curve \(\vec r(s)\) that exists in
three-dimensional space, i.e.
\begin{equation*}
	\vec r (s) = (x(s),y(s),z(s)).
\end{equation*}
The space curve \(\vec r\) is equipped with tangent \(\vec t(s) = \vec r'(s)\) in addition to a normal vector \(\vec n(s)\) and binormal vector \(\vec b(s)\). We then define the orthonormal triad \(\{\vec t, \vec n, \vec b\}\). Then the \href{https://en.wikipedia.org/wiki/Darboux_vector}{Darboux vectors} are \(\vec t ' = \vec \kappa \times \vec t \), \(\vec n ' = \vec \kappa \times \vec n\) and \(\vec b ' = \vec \kappa \times \vec b\) where \(\vec \kappa = \vec t \kappa_t + \vec n \kappa_n + \vec b \kappa_b \) with \(\kappa_t = 0\) since we assume no rotation around the neutral axis. Then we obtain
\begin{align}
	\vec t'        & = \kappa_n \vec n - \kappa_b \vec b \label{eq:t-prime-darboux} \\
	\qquad \vec n' & = -\kappa_b \vec t \label{eq:n-prime-darboux}                  \\
	\qquad \vec b' & = \kappa_n \vec t \label{eq:b-prime-darboux}
\end{align}
Using this, we can define a mapping \(\vec{R} : (s,\eta,\beta) \mapsto (x,y,z)\) for any material point in the cross section plane along \((\vec n, \vec b)\) of our beam to Cartesian space. We define this as
\begin{equation}\label{eq:material-mapping}
	\vec R (s,\eta, \beta) = \vec r (s) + \eta \: \vec n(s) + \beta \:\vec b(s).
\end{equation}
Then, for a fixed value of \((\eta, \beta)\) we have from \eqref{eq:t-prime-darboux}-\eqref{eq:material-mapping} that
\begin{align*}
	\pdif{s} \vec R & = \vec r' + \eta \: \vec n' + \beta \: \vec b'                  \\
	                & = \vec t + \eta \: \vec n' + \beta \: \vec b'                   \\
	                & = \vec t + \eta \: (-\kappa_b \vec t) + \beta (\kappa_n \vec t) \\
	\pdif{s} \vec R & = (1 - \eta \kappa_b + \beta \kappa_n) \vec t
\end{align*}
Then it follows that we have stretch ratio along the neutral axis which is
\begin{equation*}
	\lambda(s,\eta,\beta) = \norm{ \pdif{s} \vec R } = 1 - \eta \kappa_b + \beta \kappa_n.
\end{equation*}
Under the assumptions of small strains (\(\abs{\eta \kappa_b}, \abs{\beta \kappa_n} \ll 1\)), we use the engineering axial strain and
\begin{equation*}
	\varepsilon_{tt}(s,\eta,\beta) = \lambda - 1 = \varepsilon_{m}(s) -\eta \kappa_b(s) + \beta \kappa_n(s).
\end{equation*}
where \(\varepsilon_{m}(s) = 0\) under our assumption of inextensibility, e.g.
\begin{equation} \label{eq:strain-tt}
	\varepsilon_{tt}(s,\eta,\beta) = \lambda - 1 = -\eta \kappa_b(s) + \beta \kappa_n(s).
\end{equation}
We then obtain for our stresses
\begin{equation*}
	\sigma_{tt}(s, \eta, \beta) = E \varepsilon_{tt}(s,\eta,\beta)
\end{equation*}
which is equivalent to \(\tens \sigma_{tt} \: \vec t \otimes \vec t\), with \(\sigma_{tn} = \sigma_{nt} = 0\), \(\sigma_{tb} = \sigma_{bt} = 0\) and \(\sigma_{nn} = \sigma_{bb} = 0\). Then define the second moments of area
\begin{equation*}
	I_{\eta\eta} = \int_A \eta^2 \odif{A}, \quad I_{\beta\beta} = \int_{A} \beta^2 \odif{A}, \quad I_{\eta\beta} = \int_{A} \eta \beta \odif{A}, \quad A = \int_A\odif{A}
\end{equation*}
about the neutral axis in the material cross section. Define our resultant force \(N(s)\) as
\begin{equation*}
	N(s) = \int_A \tens \sigma_{tt} \odif{A} = 0,
\end{equation*}
and our second moment of inertia \(\vec M(s)\) as
\begin{equation*}
	\vec M(s) = \int_A \vec r_{\perp} \times \vec p \odif{A}, \quad \vec r_{\perp} = \eta \vec n + \beta \vec b.
\end{equation*}
With \(\vec r_{\perp} \times \vec t = - \eta \vec b + \beta \vec n\) and \(\tens \sigma_{tt} = E( - \eta \kappa_b + \beta \kappa_n)\), then
\begin{equation*}
	N(s) = EA\varepsilon_m(s),
\end{equation*}
where \(\varepsilon_m(s)\) is zero due to an assumption of no axial strain, and so
\begin{equation*}
	N(s) = 0.
\end{equation*}
In addition,
\begin{equation*}
	\vec M(s) = E \left[\left(I_{\eta\eta} \kappa_b - I_{\eta\beta} \kappa_n\right) \vec b + \left(-I_{\eta\beta} \kappa_b + I_{\beta\beta}\kappa_{n}\right)\vec n\right].
\end{equation*}
If we treat \((\eta,\beta)\) as principal bending axes, then \(I_{\eta\beta} = 0\) and this decouples to
\begin{equation*}
	M_b = E I_{\eta\eta} \kappa_b, \quad M_n = E I_{\beta\beta} \kappa_n,
\end{equation*}
which are the familiar moment-curvature laws. At a point, the strain energy is the area under the \(\sigma\)--\(\varepsilon\) curve
\begin{equation*}
	\psi = \int_0^{\varepsilon} \sigma(\tilde \varepsilon) \odif{\tilde \varepsilon}.
\end{equation*}
with Hooke's law \(\sigma = E \varepsilon\),
\begin{equation*}
	\psi = \int_0^{\varepsilon_{tt}} E \tilde \varepsilon \odif{\tilde \varepsilon} = \frac12 E \varepsilon_{tt}^2.
\end{equation*}
We may compute the energy per unit length, by the section integral
\begin{equation*}
	w(s) = \int_A \psi \odif{A} = \frac{1}{2} E \int_A \varepsilon_{tt}^2 \odif{A}.
\end{equation*}
Using \eqref{eq:strain-tt} then our energy per unit length is
\begin{equation*}
	w(s) = \frac12 E (I_{\eta\eta} \kappa_b^2 + I_{\beta\beta} \kappa_n^2 - 2I_{\eta\beta}\kappa_b \kappa_n) = \frac12 E (I_{\eta\eta}\kappa_b^2 + I_{\beta\beta}\kappa_n^2),
\end{equation*}
under the assumption of in-plane bending only, and thus setting \(I_{\eta\beta} = 0\). In addition, we assume that the beam is isotropic in such a way that we may write \(I_{\eta\eta} = I_{\beta\beta} = I\) and then we may write
\begin{equation}
	w(s) = \frac12 EI (\kappa_b^2 + \kappa_n^2) = \frac12 EI \abs{\vec t'(s)}^2  = \frac12 EI \abs{\vec r''(s)}^2.
\end{equation}

\paragraph {Small Angle Euler-Bernoulli Beam}
To complete the derivation of the classic Euler-Bernoulli beam, we may make a
small-angle or small slope argument: The exact scalar curvature \(\kappa\) for
the in-plane bending of the beam is
\begin{equation*}
	\kappa = \norm{\odv{\vec t}{s}} = \norm{ \odv{\vec t}{x} \odv{x}{s}}
\end{equation*}
\begin{equation*}
	w_y(x) = \frac{1}{2} EI (y''(x))^2
\end{equation*}
and that this is balanced by a body force \(q_y(x)\) (for \(x\in(0,L)\)) that is applied orthogonally in the principal Cartesian axis \(y\). Then the total potential energy of the beam is
\begin{equation*}
	\Pi[y] = \int_0^L \left(\frac12 EI (y''(x))^2 - q(x)y\right) \odif{x}
\end{equation*}
Take the first variation and integrate by parts twice to obtain
\begin{equation*}
	EI y^{(4)}(x) = q(x), \quad x \in (0,L).
\end{equation*}

% \subsection{Variational Problem}

% The strong form of the Euler--Bernoulli beam theory is \begin{equation*}
% \odv*[2]{\left[EI \odv[2]{y}{x}\right]}{x} = q(x) \end{equation*} Multiply the
% above by a test function \(v\in V\) and integrate over the domain \([0,L]\) of
% the beam to obtain \begin{equation*} \int_0^L \odv*[2]{\left[EI
% \odv[2]{y}{x}\right]}{x} v(x) \odif {x} = \int_0^L q(x) v(x)\odif{x}
% \end{equation*} Using integration by parts on the left hand side,
% \begin{align*} \int_0^L \odv*[2]{\left[EI \odv[2]{y}{x}\right]}{x} v(x) \odif
% {x} & = \left[\odv{}{x}\left(EI \odv[2]{y}{x}\right) v(x)\right]_0^L -
% \int_0^L \odv*{\left[EI \odv[2]{y}{x}\right]}{x} \odv{v}{x} \odif{x} \\ & =
% \left[\odv*{\left(EI \odv[2]{y}{x}\right)}{x} v(x)\right]_0^L -
% \left[EI\odv[2]{y}{x} \odv{v}{x}\right]_0^L + \int_0^L EI \odv[2]{y}{x}
% \odv[2]{v}{x} \odif{x} \\ & = \int_0^L EI \odv[2]{y}{x} \odv[2]{v}{x} \odif{x}
% + \left[\odv*{\left(EI \odv[2]{y}{x}\right)}{x} v(x)- EI\odv[2]{y}{x}
% \odv{v}{x}\right]_0^L \end{align*} Then we can rewrite the equation as
% \begin{equation*} \int_0^L EI \odv[2]{y}{x} \odv[2]{v}{x} \odif{x} = \int_0^L
% qv\odif{x} - \left[\odv*{\left(EI \odv[2]{y}{x}\right)}{x} v- EI\odv[2]{y}{x}
% \odv{v}{x}\right]_0^L, \qquad \forall v \in V. \end{equation*} Therefore, our
% variational problem is:

% \begin{weakform} Find \(y \in \overline{y} + V\) such that \begin{equation*}
% a(y,v) = \ell(v) + \ell_{\rm{tip}}(v), \quad \forall v \in V \end{equation*}
% with our bilinear form \(a\) defined as \begin{equation*} a(y,v) = \int_0^L EI
% y'' v'' \odif{x} \end{equation*} and our linear functionals \(\ell\) and
% \(\ell_{\rm{tip}}\) defined as \begin{equation*} \ell(v) = \int_0^L qv
% \odif{x} \qquad \text{and} \qquad \ell_{\rm{tip}}(v) = [(EI y''')v -
% EIy''v']_0^L. \end{equation*} \end{weakform}

% \section{Large Angle Euler-Bernoulli Beam}

% We consider the strong form \begin{equation} \odv*[2]{\left(EI(s)
% \odv[2]{\vec{w}}{s}\right)}{s} - \odv*{\left(\lambda
% \odv{\vec{w}}{s}\right)}{s} = \vec{q}(s). \end{equation} subject to an
% inextensibility constraint \begin{equation*} \norm{\odv{\vec{w}}{s}} = 1.
% \end{equation*} Multiply the equilibrium equation by \(\vec{v} \in V_0\) and
% integrate over \((0,L)\) to get \begin{equation*} \int_0^L
% \left(\left(EI\:\vec{w}^{(2)}\right)'' - (\lambda \vec{w}')'\right) \cdot
% \vec{v} \odif{s} = \int_0^L \vec{q} \cdot \vec{v} \odif{s}. \end{equation*}
% Then let \(\vec{B} \coloneqq EI \vec{w}''\) \begin{align*} \int_0^L \vec{B}''
% \cdot \vec{v} \odif{s} & = \left[\vec{B}' \cdot \vec{v}\right]_0^L - \int_0^L
% \vec{B}' \cdot \vec{v}' \odif{s} \\ & = [\vec{B}' \cdot \vec{v} - \vec{B}
% \cdot \vec{v}']_0^L + \int_0^L \vec{B} \cdot \vec{v}'' \odif{s} \\ \int_0^L
% (EI\:\vec{w}'')'' \cdot \vec{v} \odif{s} & = \int_0^L EI \vec{w}'' \cdot
% \vec{v}'' \odif{s} + \left[(EI \vec{w}'')' \cdot \vec{v} - (EI \vec{w}'')
% \cdot \vec{v}'\right]_0^L \end{align*} For the multiplier term \begin{align*}
% -\int_0^L (\lambda \vec{w}')' \cdot \vec{v} \odif{s} & = - \left[\lambda
% \vec{w} \cdot \vec{v}\right]_0^L + \int_0^L \lambda \vec{w} \cdot \vec{v}
% \odif{s}. \end{align*} To treat the inextensibility constraint, multiply the
% constraint by a test function \(\mu \in \Lambda\). Then \begin{equation*}
% \int_0^L \mu(\norm{\vec w'}^2 - 1) \odif{s} = 0. \end{equation*}

% \begin{weakform} Then our weak problem is to find \((\vec w, \lambda) \in V
% \times \Lambda\) sicj that for all \((\vec v, \mu) \in V_0 \times \Lambda\)
% \begin{equation*} a(\vec{w},\vec{v}) + b(\lambda; \vec{w},\vec{v}) =
% \ell(\vec{v}) \end{equation*} with \begin{equation*} a(\vec{w},\vec{v}) =
% \int_0^L EI \vec{w}'' \cdot \vec{v}'' \odif{s}, \qquad b(\lambda ;\vec{w},
% \vec{v}) = \int_0^L \lambda \vec{w}' \cdot \vec{v}' \odif{s} \end{equation*}
% and \begin{equation*} \ell(\vec{v}) = \int_0^L \vec{q} \cdot \vec{v} \odif{s}
% + \left[(EI \vec{w}'')' \cdot \vec{v} - (EI \vec{w}'') \cdot \vec{v}' -
% \lambda \vec{w}' \cdot \vec{v}\right]_0^L \end{equation*} subject to
% \begin{equation*} \int_0^L \mu (\norm{\vec{w'}}^2 - 1) \odif{s} = 0.
% \end{equation*} \end{weakform}
